\documentclass[10pt,a4paper]{amsart}
\usepackage[margin=19mm]{geometry}
\usepackage{amsmath,amssymb,mathtools}
\usepackage[scr=rsfs,cal=euler]{mathalfa}
\usepackage{xfrac}
\usepackage[unicode,hidelinks,bookmarks=false]{hyperref}
\newcommand{\ie}{i.e.\ }
\newcommand{\cf}{cf.\ }
\newcommand{\resp}{respectively\ }
\newcommand{\Subsection}{\subsection}
\providecommand{\nobreakdash}{\nobreak\hbox{-}}
\setlength{\emergencystretch}{2em}
\multlinegap0pt
\usepackage{amssymb}
\usepackage{enumerate}
\usepackage{cancel}
\newtheorem{proposition}{Proposition}
\newtheorem{lemma}{Lemma}
\title{Optimal bilinear control restricted to the three-dimensional chemo-repulsion model with potential production}
\author{Francisco Guillen-Gonzalez, Exequiel Mallea-Zepeda, Maria A. Rodriguez-Bellido, Elder J. Villamizar-Roa}
\date{Source: arXiv:2603.18900v1 (CC BY 4.0)}
\begin{document}
\maketitle

\noindent\emph{Source and licence.} Optimal bilinear control restricted to the three-dimensional chemo-repulsion model with potential production. Version arXiv:2603.18900v1 (CC BY 4.0), distributed by arXiv under the Creative Commons Attribution 4.0 International licence, http://creativecommons.org/licenses/by/4.0/, as recorded in the arXiv licence metadata for this exact version and shown on the abstract page https://arxiv.org/abs/2603.18900v1. Full licence notice: \texttt{licenses/arxiv-2603.18900-CC-BY-4.0.txt}.\par\smallskip

\noindent\textit{Editorial scope.} Counted unit: the article's lemma der-oper (source0.tex lines 1874--1897). The article's complete original proof of it is delivered with it (source0.tex lines 1898--2075, one \texttt{proof} environment of 178 lines). No mathematics is written, repaired or re-derived: the statement and the proof are the article's own lines, in the article's own notation, and no retained line is substituted or re-flowed.  Retained uncounted as the source-local context the proof needs: lines 450--458.  Nothing was omitted from any retained span, so the package carries no omission record.  No retained line contains a citation, so the wrapper carries no bibliography and no bibliography entry.  No retained line contains a figure environment, an image-inclusion command or drawing code, so no image is delivered and the package declares no asset dependency.  Editorial formatting only, in addition: an AMS \texttt{amsart} preamble that restates the article's own declarations and macros used by the retained lines, namely the environments \texttt{proposition}, \texttt{lemma} on separate counters, together with the standard packages those lines need.  The retained environments print in this document as Proposition 1, Lemma 1 in order of appearance, whereas the article prints the same items with the numbering of its own full document.  The complete native source of the article as distributed in the arXiv e-print for this exact version is bundled unchanged as \texttt{sources/196735/source0.tex}.
\begin{proposition}\label{ident}
Let $F\in Z'.$ Then, there exists 
$$(g^0,{\bf g}^1)\in L^{5\alpha(p)/(2\alpha(p)+3)}(Q)\times (L^{5\alpha(p)/(\alpha(p)+3)}(Q))^3$$ such that
$$\langle F,\varphi\rangle=
\int_0^T\int_\Omega g^0\varphi+\int_0^T\int_\Omega {\bf g}^1\cdot\nabla \varphi,
\quad \forall\varphi\in Z.
$$
In fact, $F=g^0-\nabla\cdot {\bf g}^1$ in $Z'$.
\end{proposition}

\begin{lemma}\label{der-oper}
The operator $\mathcal{S},$ from ${X}_{u}  \times {X}_{v} \times L^{5/2}(L^{5/2+}(\Omega_c))$ to $Y_u \times Y_v,$ is continuously Fr\'echet-differentiable and its derivative at an arbitrary point 
$z^*=(u^*,v^*,f^*)\in {X}_{ u} \times {X}_{v} \times L^{5/2}(L^{5/2+}(\Omega_c)),$ in the direction $z=(U,V,F)\in {X}_{ u} 
\times{X}_{v}\times L^{5/2}(L^{5/2+}(\Omega_c)) $ is the linear and bounded operator $\mathcal{S}'(z^*)=(\mathcal{S}_1'(z^*),\mathcal{S}_2'(z^*)): {X}_{ u} 
\times{X}_{v}\times L^{5/2}(L^{5/2+}(\Omega_c)) \to Y_u \times Y_v$ defined by
\begin{eqnarray}\label{sss-2-u}
&&\langle\mathcal{S}_1'(z^*)[z],\overline{u}\rangle_{Z'}
=
\int_0^T\langle\partial_tU,\overline{u}\rangle_{Z'}
+\int_0^T\int_\Omega\nabla U\cdot\nabla\overline{u}
%+\int_0^T\int_\Omega\nabla U\cdot\nabla\overline{u}
\nonumber\\
&&+\int_0^T\int_\Omega\left(u^*\nabla V+U\nabla v^*\right)\cdot\nabla\overline{u}
%\nonumber\\
-\int_0^T\int_\Omega\left(rU - \mu \,p \, |u^*|^{p-2} u^* U\right)\overline{u},
\quad \forall\,\overline{u}\in Z,
\end{eqnarray}
and 
\begin{equation} \label{sss-2-v}
\mathcal{S}_2'(z^*)[z]=
\partial_tV-\Delta V+V-p \, |u^*|^{p-2} u^* U-f^*V\,1_{\Omega_c}-Fv^*\,1_{\Omega_c}
\quad \hbox{in $L^{5/2}(Q)$}.
\end{equation}
\end{lemma}

\begin{proof}
% Let $s^*=[u^*,v^*,f^*]$ and $s=[U,V,F]$ two elements belonging to 
%$\widehat{X_u} \times\widehat{X}_{5/2}\times { L_t^{5/2}L_{\bf x}^{5/2+}(Q_c)}$. 
%
%\
We first analyze the Fr\'echet differentiability of $\mathcal{S}_1$.  Let $z^*=(u^*,v^*,f^*)$ and $z=(U,V,F)$ two elements of
${X}_u\times {X}_v\times L^{5/2}(L^{5/2+}(\Omega_c))$. We have
\begin{equation}\label{sss-3}
\left\{
\begin{array}{rcl}
\mathcal{S}_1(z^*+z)-\mathcal{S}_1(z^*)
&=&\partial_tU-\Delta U-\nabla\cdot(u^*\nabla V)-\nabla\cdot(U\nabla v^*)-\nabla\cdot(U\nabla V)\\
&& -rU + \mu\Big( |u^*+U|^p - |u^*|^p \Big) \ \ \mbox{ in } Y_u=Z'.
\end{array}
\right.
\end{equation}

Using that the real function  $\varphi(y)= |y|^p$ is continuously-differentiable in $\mathbb{R}$ (recall that $p>1$), from the mean-value theorem  there exists a measurable function $\theta:Q\rightarrow \mathbb{R}$ attaining intermediate values between the ones of $u^*$ and $u^*+U$ such that
\begin{equation}\label{sss-4}
|u^*(t,x)+U(t,x)|^p-|u^*(t,x)|^p=  p \, |\theta(t,x)|^{p-2}\theta(t,x) U(t,x),\ a.e.\ (t,x)\in Q.
\end{equation}
%
Furthermore, by the dominated convergence theorem, we have 
%

\begin{equation}\label{sss-4-1}
|\theta|^{p-2}\theta\to |u^*|^{p-2} u^* \hbox{ in $L^{\beta(p)/(p-1)}(Q)$, }
%for any $q\le 5/(p-1)$}
\mbox{ as } U\to 0 \hbox{ in $L^{\beta(p)}(Q)$}.
\end{equation}

Then, equality \eqref{sss-2-u} together with \eqref{sss-3}
 and \eqref{sss-4}  and Proposition \ref{ident}
imply

\begin{equation}\label{sss-5}
\begin{array}{rcl}
&&\|\mathcal{S}_1(z^*+z)-\mathcal{S}_1(z^*)-\mathcal{S}_1'(z^*)[z]\|_{Y_u}
\\
 && \quad \le  \|\nabla\cdot(U\nabla V)\|_{Y_u}
 +\mu p\,\|( |u^*|^{p-2} u^* - |\theta|^{p-2}\theta )U\|_{Y_u}\\
 &&\quad 
 \le \|U\nabla V\|_{L^{5\alpha(p)/(\alpha(p)+3)}(Q)}
 +\mu p\,\|( |u^*|^{p-2} u^* - |\theta|^{p-2}\theta )U\|_{L^{5\alpha(p)/(2\alpha(p)+3)}(Q)}.
\end{array}
\end{equation}

 Now, since $(U, V)\in X_u\times X_v$, then in particular $(U,\nabla V)\in L^{\beta(p)}(Q)\times L^5(Q)$. 
 Moreover, taking into account that $(5/3)\alpha(p)=\beta(p)$ for any $p>1$, 
 from the H\"older inequality we obtain
\begin{eqnarray}
\|U\nabla V \|_{ L^{5\alpha(p)/(\alpha(p)+3)}(Q)}
&\leq& 
\|U\|_{L^{5\alpha(p)/3}(Q)}\|\nabla V\|_{L^{5}(Q)}
%\nonumber\\
\le C\|U\|_{{X}_{ u}}\|V\|_{{X}_{ v}}.
%\le 
%C \|[U,V,F]\|^2_{{X}_{ u}\times{X}_{ v}\times{X}_f}.
\label{sss-6}
\end{eqnarray}


On the other hand, 
%since $5r(p)/(2r(p)+3)\le 5/2$ and $5s(p)/(2s(p)-5)\le s(p)/(p-1)$ for any $p>1$, 
from the H\"older inequality we have
%
\begin{eqnarray}
&&
\|( |u^*|^{p-2} u^* - |\theta|^{p-2}\theta )U\|_{L^{5\alpha(p)/(2\alpha(p)+3)}(Q)}
 \nonumber \\
&& \quad \le \|( |u^*|^{p-2} u^* - |\theta|^{p-2}\theta )\|_{L^{5/2}(Q)} 
 \|U\|_{L^{5\alpha(p)/3}(Q)}
% \nonumber \\
% &&  \quad \le   \,  C \,  \| |u^*|^{s-2} u^* - |\theta|^{s-2}\theta \|_{L^{\beta(p,s)/(s-1)}}\|U\|_{L^{\beta(p,s)}}
\nonumber \\
 &&\quad \le   \,  C \,  \| |u^*|^{p-2} u^* - |\theta|^{p-2}\theta \|_{L^{5/2}(Q)}\|U\|_{X_u}.
\label{sss-6-1}
\end{eqnarray}

%
Therefore, from \eqref{sss-5}-\eqref{sss-6-1}  
and taking into account \eqref{sss-4-1}, we deduce that
\begin{eqnarray*}
&&\lim_{\| z \|_{{X}_{u}\times {X}_{v}
\times L^{5/2}(L^{5/2+})} \to0}
\frac{\|\mathcal{S}_1(z^*+z)-\mathcal{S}_1(z^*)-\mathcal{S}_1'(z^*)[z]\|_{Y_u}}
{\|z \|_{ {X}_{ u}\times {X}_{v}
\times L^{5/2}(L^{5/2+})}}\\
&&
 \le \lim_{\| z \|_{{X}_{ u}\times{X}_{v}\times L^{5/2}(L^{5/2+})} \to0}
C \left(\| z \|_{{X}_{u}\times{X}_{v}
\!\times L^{5/2}(L^{5/2+})}
+
\!\| |u^*|^{p-2} u^* - |\theta|^{p-2}\theta \|_{L^{ 5/2}(Q)}\right)=0.
\end{eqnarray*}


Consequently, the operator $\mathcal{S}_1$ is Fr\'echet-differentiable and its Fr\'echet derivative is given in
\eqref{sss-2-u}. The continuity of the operator 
$$z^*\in X_u\times X_v\times L^{5/2}(L^{5/2+}(\Omega_c)) \longmapsto S_1'(z^*)[\cdot]\in  \mathcal{L}(X_u\times X_v\times L^{5/2}(L^{5/2+}(\Omega_c));Y_u)$$ can be deduced using the same type of estimates developed previously. In fact, one has

$$
\| \mathcal{S}_1'(z_1^*)-\mathcal{S}_1'(z_2^*)\|_{ \mathcal{L} }%(X_u\times X_v\times L^{5/2}(L^{5/2+}(\Omega_c));Y_u)} 
\le C(\| u_1^*-u_2^*\|_{L^{\beta(p)}(Q)} + \| \nabla (v_1^*-v_2^*) \|_{L ^{5}(Q)})
+ \| |u_1^*|^{p-2}u_1^*-|u_2^*|^{p-2}u_2^* \|_{L^{5/2}(Q)}).
$$


Now, we will analyze the Fr\'echet differentiability of $\mathcal{S}_2$. Notice that
%\begin{equation}\label{sss-3-2}
$$
\mathcal{S}_2(z^*+z)-\mathcal{S}_2(z^*)
=\partial_t V-\Delta V+ V -
\vert u^*+U \vert^{p} 
+|u^*|^p - (f^* V\,+ F v^* + F V)\,1_{\Omega_c}.
%\end{equation}
$$
Then, 
%defining
%$$
%\mathcal{S}_2'(s^*)[s]=\partial_tV-\Delta V+V- p(u^*)^{p-1}U
%- f V\,1_{\Omega_c} - F v \,1_{\Omega_c},
%$$
from \eqref{sss-2-v} and \eqref{sss-4} we get
%
\begin{eqnarray*}
&&\|\mathcal{S}_2(z^*+z)-\mathcal{S}_2(z^*)-\mathcal{S}_2'(z^*)[z]\|_{L^{ 5/2}(Q)}\\
&&\ \ \le 
p\, \| ( |u^*|^{p-2} u^* -|\theta|^{p-2}\theta)U\|_{L^{ 5/2}(Q)} + \| F \, V\|_{L^{ 5/2}(Q)}\\
&&\ \ \le  \||u^*|^{p-2} u^* -|\theta|^{p-2}\theta\|_{L^{5\beta(p)/(2\beta(p)-5)}(Q)}
\|U\|_{L^{ \beta(p)}(Q)}\\
&&\ \ \ +   \| F \|_{L^{5/2}(L^{5/2+}(\Omega_c))}\|V\|_{L^\infty(L{^{\infty-}})}\\
%&&\ \ \le C  \||u^*|^{p-2} u^* -|\theta|^{p-2}\theta\|_{L^{\beta(p,s)/(p-1)}(Q)}\|U\|_{L^{\beta(p,s)}(Q)}
%\ \ +   C\| F \|_{L^{5/2}(L^{5/2+}(\Omega_c))}\|V\|_{X_v}\\
&&\ \ \le C  \||u^*|^{p-2} u^* -|\theta|^{p-2}\theta\|_{L^{5\beta(p)/(2\beta(p)-5)}(Q)}\|U\|_{X_u}+ C\| F \|_{L^{5/2}(L^{5/2+}(\Omega_c))}\|V\|_{X_v}.
\end{eqnarray*}

%By applying the H\"older inequality we have
%\begin{eqnarray}
%\|( |u^*|^{p-2} u^* - |\theta|^{p-2}\theta)U\|_{L^{5/2}(Q)}
%&\le&  \||u^*|^{p-2} u^* -|\theta|^{p-2}\theta\|_{L^{5}}\|U\|_{L^{ {5}}(Q)}\label{sss-6-1-2}
%\end{eqnarray}
%Similarly, we can obtain
%$$
%\mu q\|(u^*+U)^{q-1}U\|_{L^5(Q)}\le C\|[U,V,F]\|^2_{\widehat{X}_5\times\widehat{X}_5\times L^5(Q_c)}.
%$$

Therefore,  taking into account \eqref{sss-4-1}  and that $\frac{5\beta(p)}{2\beta(p)-5}\le \frac{\beta(p)}{p-1}$
(in fact, $\frac{5\beta(p)}{2\beta(p)-5}=5$ if $p\le2$ and $\frac{5\beta(p)}{2\beta(p)-5}=\frac{5(p-1)}{2p-3}<5$ if $p>2$),
we deduce
\begin{eqnarray*}
&&\lim_{\|z \|_{{X}_{ u}\times{X}_{v}
\times L^{5/2}(L^{5/2+})} \to 0}
\frac{\|\mathcal{S}_2(z^*+z)-\mathcal{S}_2(z^*)-\mathcal{S}_2'(z^*)[z]\|_{L^{ 5/2}(Q)}}
{\|z \|_{{X}_{u}\times{X}_{v} \times L^{5/2}(L^{5/2+})}}\\
&&
\quad \le \lim_{\| z \|_{{X}_{u}\times{X}_{v} \times L^{5/2}(L^{5/2+})} \to0}
C \left( \| |u^*|^{p-2} u^* -|\theta|^{p-2}\theta\|_{L^{5\beta(p)/(2\beta(p)-5)}(Q)} +  \| F \|_{L^{5/2}(L^{5/2+})} \right)=0.
\end{eqnarray*}

%{\color{blue}
%DEBE SER $2s(p)/(s(p)-2) \le 5s(p)/(2s(p)-5)\le s(p)/(p-1)$. ES IGUAL PARA $p=2$ Y ES ESTRICTO PARA $p<2$ 
%Y TAMBIEN ES ESTRICTO PARA $p>2$
%}

Consequently, the operator $\mathcal{S}_2$ is Fr\'echet-differentiable and its Fr\'echet derivative is given in 
%$S_2'(s^*)[s]$
\eqref{sss-2-v}. Again, the continuity of the operator 
$$z^*\in X_u\times X_v\times L^{5/2}(L^{5/2+}(\Omega_c)) \longmapsto S_2'(z^*)[\cdot]\in  \mathcal{L}(X_u\times X_v\times L^{5/2}(L^{5/2+}(\Omega_c)); Y_v)
$$
 can be proved. In fact, we get 
\begin{eqnarray*}
\| \mathcal{S}_2'(z_1^*) - \mathcal{S}_2'(z_2^*)\|_{\mathcal{L}} %(X_u\times X_v\times L^{5/2}(L^{5/2+}(\Omega_c));Y_v)}
 &\le& C( \| |u_1^*|^{p-2}u_1^*-|u_2^*|^{p-2}u_2^* \|_{L ^{5\beta(p)/(2\beta(p)-5)}(Q)} \\
 &&+ \| v_1^*-v_2^* \|_{L^{\infty}(L^{\infty-})} + \| f_1^*-f_2^* \|_{L^{5/2}(L^{5/2+}(\Omega_c))}).
\end{eqnarray*}
%
\end{proof}

\end{document}
